\documentclass[11pt,a4paper]{article}

\usepackage[a4paper, left=1in, right=1in, top=1in, bottom=1in]{geometry}
\usepackage[hidelinks]{hyperref}
\usepackage{amsmath, amssymb}
\usepackage{graphicx}
\usepackage{booktabs}
\usepackage{parskip}
\usepackage{xcolor}
\usepackage{mdframed}
\usepackage{enumitem}
\usepackage{indentfirst}

\setlength{\parindent}{0pt}
\setlength{\parskip}{8pt}

\definecolor{noteblue}{RGB}{220,235,250}
\definecolor{noteyellow}{RGB}{255,250,220}
\definecolor{notegreen}{RGB}{220,245,225}

\newmdenv[backgroundcolor=noteblue, linewidth=0pt, innerleftmargin=10pt,
          innerrightmargin=10pt, innertopmargin=8pt, innerbottommargin=8pt]{keybox}

\newmdenv[backgroundcolor=noteyellow, linewidth=0pt, innerleftmargin=10pt,
          innerrightmargin=10pt, innertopmargin=8pt, innerbottommargin=8pt]{intuitionbox}

\newmdenv[backgroundcolor=notegreen, linewidth=0pt, innerleftmargin=10pt,
          innerrightmargin=10pt, innertopmargin=8pt, innerbottommargin=8pt]{resultbox}

\begin{document}

\begin{center}
{\LARGE \textbf{How My Thesis Works: A Self-Study Guide}}\\[0.5em]
{\large From Classical Metrics to Functional Forecasting}\\[0.5em]
{\normalsize Zeynep Bozoklu \quad --- \quad Personal learning notes}
\end{center}

\vspace{1em}

These notes explain every concept in my thesis starting from what I already
know from first-year econometrics. The goal is to be able to explain any
part of the work clearly and confidently in a seminar or presentation.
The actual numbers from the results appear throughout so the explanations
stay grounded.

\tableofcontents

\newpage

% ============================================================
\section{Why This Problem Matters}
% ============================================================

Turkey's electricity grid operator needs to forecast how much electricity
will be consumed in each of the 24 hours of the following day. This is
called a \textbf{day-ahead load forecast}. Grid operators use these
forecasts to schedule power plants, manage reserves, and price electricity
in the wholesale market.

The Turkish operator publishes an official forecast. My research question
is simple: \textit{can a statistical model do better?}

The answer, for the model I build: yes, by about \textbf{19\% in mean
absolute error} and \textbf{19\% in mean absolute percentage error}
relative to the official forecast.

Why might a statistical model beat the official forecast? The official
forecast is itself the output of a complex proprietary system. But a
statistical model trained on historical patterns can sometimes learn
systematic biases or regularities that the official model misses.

% ============================================================
\section{The Economic Setting}
% ============================================================

\subsection{Why the load curve has the shape it has}

The 24-hour electricity load curve is a direct imprint of human behavior
and economic activity. Understanding its shape is essential for explaining
why the forecasting problem is hard and why the functional data approach
is natural.

\textbf{Overnight trough (hours 0--5).}
Most people are asleep. Factories running night shifts are the main
consumers, along with street lighting and always-on infrastructure
(servers, refrigeration, pumping stations). Demand is low and stable,
varying mainly by season (heating in winter, no AC). This is the easiest
part of the curve to forecast --- weekly seasonality (lag 7) captures most
of the variation. Our model gains 40--60\% in MAPE over the official
forecast at these hours.

\textbf{Morning ramp-up (hours 6--8).}
The single most economically consequential transition in the curve.
Within two hours, electricity demand roughly doubles. The drivers are:
residential consumers waking up (lighting, coffee machines, showers),
commercial buildings switching on HVAC, and industrial facilities
starting morning shifts. The \emph{timing} of this ramp depends
critically on temperature: on a hot summer day, air conditioning turns
on immediately at sunrise; on a mild day, the ramp is gentler.
This is the one part of the daily curve where our model loses to the
official forecast (by up to 24\% MAPE at hour 8), almost certainly
because the official forecast incorporates weather forecasts that
we do not have access to.

\textbf{Midday plateau (hours 9--16).}
Industrial and commercial demand is at full intensity.
This part of the curve is dominated by large consumers ---
manufacturing, mining, construction --- whose behavior is
predictable from the day of week and seasonal calendar.
Our model wins comfortably here (20--35\% MAPE gain),
suggesting these systematic patterns are well captured by
holiday indicators, weekend flags, and lagged load-shape features.

\textbf{Evening surge (hours 17--19).}
The second major transition. Industrial activity begins to wind
down while residential demand surges: people come home, turn on
televisions, cook dinner, and in summer run air conditioning at
maximum. In Turkey this peak coincides with the hottest part of
the evening in summer and with heating demand in winter. The official
forecast is worst here: it underpredicts by 614--658 MWh on average
at hours 17--19, substantially larger than at any other time of day.
Our model still loses at hours 18--19 but with a much smaller bias.

\textbf{Night wind-down (hours 20--23).}
Residential demand falls as people go to sleep. Industrial activity
has mostly stopped. The curve declines smoothly toward the overnight
trough. Our model wins by 17--20\% MAPE here.

\begin{intuitionbox}
\textbf{Why this matters for FDA:} the daily load curve is not 24
independent numbers --- it is a single coherent trajectory shaped
by overlapping behavioral rhythms. The morning ramp, midday plateau,
and evening peak are features of the \emph{whole curve}, not of
individual hours. Modeling the curve as a functional object rather
than a vector of 24 hourly values lets us directly work with these
shape features (via FPCA components and derivative summaries) instead
of treating each hour in isolation.
\end{intuitionbox}

\subsection{Turkey's electricity market: who uses these forecasts}

Turkey liberalised its electricity market through a series of reforms
from 2001 to 2013. The key institutions are:

\textbf{TEİAŞ} (Türkiye Elektrik İletim A.Ş.) --- the transmission
system operator (TSO). It owns and operates the high-voltage grid and is
responsible for system balance: at every moment, generation must equal
consumption. TEİAŞ publishes the official day-ahead load forecast, which
is the benchmark in this thesis.

\textbf{EPIAS} (Enerji Piyasaları İşletme A.Ş.) --- the energy exchange
that operates the day-ahead market (DAM) and intraday market. The DAM
clears at noon the day before delivery, with generators and retailers
submitting hourly bids and offers.

\textbf{Generators} submit supply bids for each hour of the following
day. The clearing price depends on the demand they expect: if demand is
forecast to be high in the evening, peaking plants will be called on,
driving up prices. Generators with better demand forecasts can time their
bids more strategically --- bidding higher when they know demand will be
tight, and backing off when it will be slack.

\textbf{Retailers and large industrial consumers} submit demand bids and
purchase power at DAM prices. A large factory with a better load
forecast of its own consumption faces smaller imbalance penalties in the
real-time balancing market. Inaccurate internal forecasts expose them to
buying expensive balancing power.

\textbf{The TSO's role in scheduling:} TEİAŞ uses its load forecast to
run \emph{unit commitment} --- deciding which power plants to start up
and keep running. Starting a thermal plant (coal, gas) takes hours and
has significant startup costs. A 19\% better forecast means TEİAŞ can
commit capacity more efficiently: fewer unnecessary startups and fewer
costly real-time reserve activations.

\begin{keybox}
\textbf{The economic value of better forecasting.} A 1\% improvement
in MAPE is estimated to reduce balancing costs by millions of euros
per year in medium-sized European electricity markets (Bunn \&
Karakatsani, 2003). Turkey's electricity consumption is around
350 TWh/year and growing. Our 18.7\% MAPE improvement, if used
operationally, could represent substantial savings in reserve procurement
and balancing costs --- the exact amount would require knowledge of
Turkey's balancing market prices, but the order of magnitude is
economically meaningful.
\end{keybox}

\subsection{The systematic underprediction puzzle}

The official forecast underpredicts actual consumption at every single
hour of the day, every year in the evaluation period. The numbers:

\begin{center}
\begin{tabular}{lrr}
\toprule
Year & Official bias (MWh) & Main model bias (MWh) \\
\midrule
2023 & $+344$ & $+176$ \\
2024 & $+585$ & $+34$  \\
2025 & $+472$ & $+98$  \\
\midrule
Average & $+467$ & $+103$ \\
\bottomrule
\end{tabular}
\end{center}

\emph{Positive bias means actual $>$ forecast, i.e., the forecast
is systematically too low.}

Three observations stand out:

\begin{enumerate}
\item The underprediction worsens in 2024 (+585 MWh) before
  partially recovering in 2025. This coincides with an exceptionally
  hot summer in Turkey in 2024 driving unusually high air conditioning
  demand --- a shock the official model apparently did not absorb.

\item The underprediction is largest at evening hours (17--20),
  with a bias of 614--658 MWh. This is the residential demand surge
  that is hardest to predict from institutional/commercial patterns alone.

\item Our model's bias is much smaller (34 MWh in 2024, essentially
  zero) because the rolling-origin design continuously re-trains on the
  most recent data. As Turkish demand grew in 2024, our model adapted
  automatically.
\end{enumerate}

\subsection{Economic incentives and agent behavior}

Why would the official forecaster persistently underpredict?
Several mechanisms from the economics of forecasting apply here.

\textbf{Asymmetric loss functions.}
If the TSO faces asymmetric costs --- over-procurement of capacity
(scheduling too many plants and leaving them idle) is more politically
or contractually costly than under-procurement (buying expensive
real-time reserves) --- then rational forecast optimisation under
asymmetric loss produces a downward-biased forecast. This is a
well-known result: when the loss from over-forecasting exceeds the
loss from under-forecasting, the optimal point forecast is below the
conditional mean.

\textbf{Incentive to understate demand growth.}
A TSO that continuously reports ``demand exceeded forecast'' may face
less scrutiny than one that ``wasted money on excess capacity.''
Politically, underpredicting and then being surprised by demand
growth can be framed as evidence of a strong economy. Overpredicting
and having idle power plants on contract is visible and costly in a
different way.

\textbf{Model staleness.}
Official forecasting models at large institutions are updated
infrequently --- they require regulatory approval, internal audits,
and systems integration. A model calibrated on 2015--2019 consumption
data will systematically underpredict 2023--2025 demand if Turkish
electricity use has been structurally growing due to industrial
expansion, rising living standards, and increased AC adoption.
Our rolling-origin estimator automatically incorporates the most
recent data at each step, which is why our 2024 bias is only 34 MWh
while the official is 585 MWh.

\textbf{Ramadan and religious holiday complexity.}
Ramadan shifts the entire daily demand pattern: iftar (breaking fast
at sunset) creates a massive spike at what would otherwise be a
declining evening demand, and the overnight hours see much higher
consumption as people are awake eating, socialising, and praying.
The timing of Ramadan shifts by roughly 10--11 days per year
relative to the Gregorian calendar. An official model that does not
correctly track Ramadan will generate systematic errors that
persist for the entire month. Our model includes explicit Ramadan
and Bayram indicators constructed year by year.

\textbf{Generator strategic behavior.}
Beyond the TSO, generators themselves have incentives to shade
their demand forecasts when submitting market bids. If a generator
believes demand will be high (tight supply), bidding a high price
is rational. But if all generators believe this and bid high, the
clearing price spikes. There is a strategic complementarity:
generators want to know what others forecast, not just what actual
demand will be. The TSO's official forecast acts as a coordination
device --- generators shade their bids relative to it, which may
create feedback dynamics in the market that make the official
forecast act as a focal point rather than an accurate estimate.

\begin{intuitionbox}
\textbf{For the thesis presentation:} when an audience member asks
``why does the official forecast underpredict?'', the honest answer
is: we do not know for certain, but the data is consistent with
(a) model staleness in a period of growing demand, (b) asymmetric
institutional loss functions favouring under-procurement, and (c)
the official model not fully capturing Ramadan dynamics. The growing
gap in 2024 is most naturally explained by the model not adapting
to a structural demand increase compounded by an unusually hot summer.
\end{intuitionbox}

% ============================================================
\section{What Is Different About This Data}
% ============================================================

The raw data is hourly electricity consumption in MWh for each day from
2019 to 2025. The natural object is a \textbf{daily load curve}: for
each day $d$, we observe 24 numbers $Y_d(0), Y_d(1), \ldots, Y_d(23)$.

\begin{intuitionbox}
\textbf{The key observation:} these 24 numbers are not 24 independent
measurements. They are all part of the same daily demand curve ---
overnight is low, mornings rise steeply, afternoons plateau, evenings
surge again. Treating them as 24 separate forecasting problems throws
away all of this shape information.
\end{intuitionbox}

The shape itself is economically meaningful. Two days can have the same
\emph{average} consumption but very different \emph{shapes}: one might
have a sharp morning peak while another has a flat profile. An air
conditioning shock raises the afternoon level but not the overnight level.
A Ramadan period flattens the morning ramp and shifts demand to night.

\textbf{Functional Data Analysis (FDA)} is the branch of statistics that
treats each day's 24-hour sequence as a single mathematical object --- a
\emph{curve} --- rather than a vector of 24 numbers. This preserves the
shape information.

% ============================================================
\section{Representing the Curve: B-splines}
% ============================================================

Before doing anything statistical, we need to represent the daily curve
mathematically. We use a \textbf{B-spline basis}.

Think of it this way. A polynomial approximation says: represent any
smooth function $f(h)$ as $\sum_{m=0}^{M} c_m h^m$. This works but
has numerical problems with high-degree polynomials.

B-splines are a better version: piecewise polynomials that are smooth
at the joins. We choose $M = 12$ basis functions $B_1(h), \ldots, B_{12}(h)$
over the 24-hour domain. Then:
\[
Y_d(h) \approx \sum_{m=1}^{12} c_{dm} B_m(h)
\]
where the 12 coefficients $c_{d1}, \ldots, c_{d,12}$ are chosen to fit
the 24 hourly observations for day $d$. The basis functions $B_m(h)$ are
fixed; what changes across days are the coefficients.

\begin{keybox}
\textbf{What this buys us:} instead of working with a jagged sequence of
24 numbers, we work with a smooth curve. We can take its derivative
(useful later), integrate it, and apply the tools of functional analysis.
The 12 basis coefficients represent the curve accurately because a
smooth daily load curve does not need more than 12 degrees of freedom.
\end{keybox}

% ============================================================
\section{Functional PCA: Compressing the Curves}
% ============================================================

\subsection{Starting from classical PCA}

You likely know classical PCA. Suppose you have $n$ observations of a
$p$-dimensional vector $x_i \in \mathbb{R}^p$. PCA finds the directions
of maximum variance: the first principal component $\phi_1$ is the
direction that explains the most variance in $\{x_i\}$. The score
$\xi_{i1} = x_i' \phi_1$ is how far observation $i$ projects onto $\phi_1$.

The decomposition is:
\[
x_i = \mu + \xi_{i1}\phi_1 + \xi_{i2}\phi_2 + \cdots + \xi_{ip}\phi_p
\]
If the first $K$ components explain most of the variance, we can
approximate $x_i \approx \mu + \sum_{k=1}^K \xi_{ik}\phi_k$.

\subsection{The functional version}

\textbf{Functional PCA (FPCA)} does exactly the same thing but the
``observations'' are curves $Y_d(\cdot)$ instead of vectors. The
``directions'' are functions $\phi_k(\cdot)$ called
\textbf{eigenfunctions} (or functional principal components). The
``scores'' $\xi_{dk}$ are scalars, exactly as in classical PCA.

The decomposition is:
\[
Y_d(h) = \mu(h) + \sum_{k=1}^{K} \xi_{dk}\,\phi_k(h) + \varepsilon_d(h)
\]
where $\mu(h)$ is the mean curve, $\phi_k(h)$ is the $k$-th
eigenfunction (a 24-hour function), and $\xi_{dk}$ is the score for
day $d$ on component $k$.

\begin{intuitionbox}
\textbf{What do the eigenfunctions look like in this data?}

\textbf{FPC 1} (explains 93.4\% of variance): essentially the
``level'' component. Its shape is nearly flat across all 24 hours.
A high $\xi_{d1}$ means day $d$ has uniformly high consumption.
A low $\xi_{d1}$ means uniformly low. Think of it as ``how much
electricity was used today overall.''

\textbf{FPC 2} (explains 3.9\% of variance): the ``shape contrast''
component. Its eigenfunction is positive in the morning and evening
and negative or flat at midday. A high $\xi_{d2}$ means the
morning/evening peaks are pronounced relative to the midday level.
A low $\xi_{d2}$ means a flatter profile (more typical of summer
when air conditioning smooths the curve).
\end{intuitionbox}

Together, two components explain 97.3\% of the total variation across
all 1,096 daily curves. This is the key compression: instead of
forecasting 24 numbers, we need to forecast only 2 numbers per day ---
the two FPCA scores $(\xi_{d1}, \xi_{d2})$.

\subsection{Why K = 2?}

The choice $K=2$ follows the standard ``95\% cumulative variance'' rule.
Components 3 and 4 explain only 1.1\% and 1.0\% respectively. Including
them would add parameters to the forecasting model without meaningfully
improving the functional representation. Parsimony wins.

% ============================================================
\section{Forecasting the Scores: VAR(1,7)}
% ============================================================

\subsection{The forecasting problem after FPCA}

After FPCA, the forecasting problem becomes: given the score history
$\{\xi_{d'k} : d' < d\}$ and predetermined features $z_d$, predict
$(\xi_{d1}, \xi_{d2})$ for tomorrow.

This is a standard multivariate time series problem.

\subsection{Vector Autoregression}

A \textbf{VAR} (Vector AutoRegression) is the multivariate extension of
an AR model. If you know AR(1): $y_t = a + \beta y_{t-1} + u_t$.
A VAR(1) with $K=2$ variables is:
\[
\begin{pmatrix} \xi_{d1} \\ \xi_{d2} \end{pmatrix}
= a + A_1 \begin{pmatrix} \xi_{d-1,1} \\ \xi_{d-1,2} \end{pmatrix} + u_d
\]
where $A_1$ is a $2 \times 2$ matrix of coefficients. The VAR allows
score 1 to depend on yesterday's score 2, and vice versa.

\subsection{Why lags 1 and 7?}

Electricity demand has a strong \textbf{weekly seasonality}: Monday
demand resembles last Monday's, weekends resemble previous weekends.
So we include both lag 1 (yesterday) and lag 7 (same weekday last week):
\[
\xi_d = a + A_1\xi_{d-1} + A_7\xi_{d-7} + u_d
\]
This is the FPCA VAR(1,7) --- the core dynamic model. It is
estimated by OLS equation by equation, which is efficient when the
same regressors appear in each equation.

\subsection{Adding external features}

The VAR alone explains the \emph{dynamics} of scores but ignores
known structure: weekends have different demand than weekdays,
public holidays cause sharp drops, Ramadan shifts demand to night,
summer school breaks change the daytime pattern.

We add a vector $z_d$ of predetermined features:
\[
\xi_d = a + A_1\xi_{d-1} + A_7\xi_{d-7} + Bz_d + u_d
\]
The features $z_d$ include:
\begin{itemize}[nosep]
  \item \textbf{Calendar:} weekend indicator, day-of-year sine/cosine,
    summer school break indicator
  \item \textbf{Holiday:} fixed holidays, Ramadan and Bayram periods,
    holiday eve, post-holiday, bridge-day indicators
  \item \textbf{Load shape:} yesterday's mean, peak, range, rolling
    7-day and 28-day averages
  \item \textbf{Derivative features:} morning ramp, evening ramp,
    end-of-day drop from the \emph{smoothed derivative} of yesterday's curve
\end{itemize}

The derivative features deserve special mention. The derivative
$DY_{d-1}(h) = \frac{d}{dh}Y_{d-1}(h)$ measures how fast the
load curve was changing at each hour yesterday. A steep positive
derivative at hour 6--8 means a sharp morning ramp-up. We summarise
this by the standard deviation, maximum, minimum, and range of the
derivative, plus targeted ramps at specific transition hours. These
capture yesterday's \emph{slope behavior} in addition to its level.

\begin{keybox}
\textbf{Implementation note:} all features in $z_d$ are
\emph{predetermined} --- they are computed from data strictly before
day $d$. There is no look-ahead. This is essential for the evaluation
to be honest.
\end{keybox}

% ============================================================
\section{Rolling-Origin Evaluation}
% ============================================================

\subsection{The key principle}

Evaluating a forecasting model is different from evaluating a
regression model. In regression, we split data randomly into train
and test sets. In forecasting, the past must precede the future ---
we cannot train on 2024 data to forecast 2022.

\textbf{Rolling-origin evaluation} (also called walk-forward
cross-validation) mimics how the model would be used in practice:

\begin{enumerate}[nosep]
  \item Set a test start date: January 1, 2023.
  \item For each day $d$ from Jan 1 2023 to Dec 31 2025 (1,096 days):
    \begin{itemize}[nosep]
      \item Estimate FPCA on all data strictly before $d$
      \item Estimate the score VAR on all data strictly before $d$
      \item Produce a one-day-ahead forecast for day $d$
      \item Record the forecast error
    \end{itemize}
  \item Evaluate across all 1,096 test days.
\end{enumerate}

This design ensures \textbf{no information leakage}: the model at each
step has access only to the information that would be available in real
time. The FPCA eigenfunctions themselves are re-estimated at every step
as the training window grows.

\subsection{Expanding vs sliding window}

We use an \textbf{expanding window}: training always uses all available
past data. An alternative would be a sliding window (only the last $N$
days). Expanding windows are standard when data is non-stationary but
the long-run structure (eigenfunctions, calendar patterns) is stable.
We verified that the FPCA eigenfunctions are stable across sub-periods
--- inner products between 2019--2021 and 2022--2024 eigenfunctions
exceed 0.99 --- so the expanding window is appropriate.

% ============================================================
\section{Evaluation Metrics}
% ============================================================

\subsection{The three metrics}

For each model $m$ and each (day, hour) pair $(d,h)$, the forecast
error is $e_{dh}^{(m)} = Y_d(h) - \hat{Y}_d^{(m)}(h)$.

\[
\text{MAE}_m = \frac{1}{N}\sum_{d,h} |e_{dh}^{(m)}|
\]
\[
\text{RMSE}_m = \sqrt{\frac{1}{N}\sum_{d,h} (e_{dh}^{(m)})^2}
\]
\[
\text{MAPE}_m = \frac{1}{N}\sum_{d,h} \frac{|e_{dh}^{(m)}|}{Y_d(h)} \times 100
\]

\textbf{MAE} is the average absolute error in MWh. Directly
interpretable in the units of the problem.

\textbf{RMSE} penalises large errors more (because of the square).
If grid operators care especially about avoiding large misses on peak
days, RMSE is the more relevant metric. Note that OLS minimises RMSE
in-sample, so RMSE is also the ``coherent'' metric for OLS.

\textbf{MAPE} is percentage error. It is scale-free, which matters
when comparing across years: if Turkish electricity demand grows 5\%
between 2023 and 2025, a fixed MAE threshold becomes harder to beat
in later years. MAPE adjusts for this. In the electricity load
forecasting literature, MAPE is the most commonly reported metric.

\subsection{Relative gain vs the official forecast}

We also report the \textbf{gain} relative to the official forecast:
\[
\text{MAE gain}_m = 100 \cdot \frac{\text{MAE}_{\text{off}} - \text{MAE}_m}{\text{MAE}_{\text{off}}}
\]
Positive gain = model $m$ beats the official forecast. Negative gain =
model $m$ is worse.

\subsection{The results}

\begin{center}
\begin{tabular}{lrrrrrr}
\toprule
Model & MAE & RMSE & MAPE & MAE gain & RMSE gain & MAPE gain \\
\midrule
Seasonal naive $Y_{d-7}$ & 1912 & 3130 & 5.13\% & $-57.6\%$ & $-79.8\%$ & $-62.9\%$ \\
Official forecast         & 1214 & 1740 & 3.15\% & 0\% & 0\% & 0\% \\
FPCA VAR(1,7) scores only & 1443 & 2018 & 3.76\% & $-18.9\%$ & $-16.0\%$ & $-19.4\%$ \\
$+$ ramp/calendar         & 1238 & 1716 & 3.26\% & $-2.1\%$ & $+1.4\%$ & $-3.7\%$ \\
$+$ holiday/load regime   & 1001 & 1354 & 2.63\% & $+17.5\%$ & $+22.2\%$ & $+16.5\%$ \\
$+$ derivative features   & \textbf{980} & \textbf{1324} & \textbf{2.56\%} & $+19.2\%$ & $+23.9\%$ & $+18.7\%$ \\
$+$ residual step         & 653 & 983 & 1.70\% & $+46.2\%$ & $+43.5\%$ & $+46.0\%$ \\
\bottomrule
\end{tabular}
\end{center}

\begin{intuitionbox}
\textbf{Reading the table:} the model-building path shows each
addition. FPCA VAR(1,7) alone is \emph{worse} than the official
forecast --- the pure dynamics without any calendar information are
not enough. The ramp/calendar model barely helps. The big jump
happens when we add holiday and load-shape information: MAE drops
from 1238 to 1001. Derivative features add a further push to 980.
\end{intuitionbox}

% ============================================================
\section{Why Comparing MAE Numbers Is Not Enough: The DM Test}
% ============================================================

\subsection{The problem with just looking at numbers}

Suppose model A has MAE = 980 and model B has MAE = 990. Model A
looks better. But is this difference real, or could it arise from
random variation in the 1,096-day test period?

This is a standard statistical question: we need a hypothesis test.

\subsection{The Diebold-Mariano test}

The \textbf{Diebold-Mariano (DM) test} (Diebold \& Mariano, 1995)
answers exactly this question. Here is the logic:

Define the \textbf{loss differential} for each observation $t$:
\[
d_t = |e_t^{(\text{challenger})}| - |e_t^{(\text{official})}|
\]
If $d_t < 0$, the challenger won on observation $t$. If $d_t > 0$, the
official won. The null hypothesis is $\mathbb{E}[d_t] = 0$ (equal
expected loss). The alternative is $\mathbb{E}[d_t] < 0$ (challenger
has lower expected loss).

The DM statistic is essentially a $t$-test on $\bar{d}$, the sample
mean of the loss differentials, with a heteroskedasticity-and-autocorrelation
consistent (HAC) variance estimate. We use the Harvey, Leybourne \&
Newbold (1997) small-sample correction.

\begin{keybox}
\textbf{Connection to what you know:} the DM test is a $t$-test on
the time-series mean of $d_t$, using a long-run variance estimator
(like a Newey-West standard error in your metrics courses). It
accounts for serial correlation in the errors, which matters because
consecutive forecast errors tend to be correlated.
\end{keybox}

\subsection{The results}

\begin{center}
\begin{tabular}{lrrll}
\toprule
Challenger & $\bar{d}$ (MWh) & DM stat & $p$-value & Beats official? \\
\midrule
FPCA VAR(1,7) scores only & $+229$ & $+22.9$ & $\approx 1$ & No --- significantly worse \\
$+$ ramp/calendar         & $+25$  & $+2.6$  & $0.996$    & No \\
$+$ holiday/load regime   & $-213$ & $-24.9$ & $10^{-135}$ & Yes, $p < 0.001$ \\
$+$ derivative features   & $-233$ & $-27.4$ & $10^{-163}$ & Yes, $p < 0.001$ \\
$+$ residual step         & $-561$ & $-71.2$ & $\approx 0$ & Yes, dominant \\
\bottomrule
\end{tabular}
\end{center}

The $p$-value of $10^{-163}$ for the main model means the probability
of observing this large an improvement by chance, if both models were
truly equal, is essentially zero.

We also run the DM test \textbf{hour by hour} (24 separate tests).
The main model beats the official forecast at 19 of 24 hours, and
the difference is statistically significant (5\% level) at 18 of 24.

The 5 hours where the official forecast wins are hours 6, 7, 8
(morning ramp-up) and 18, 19 (evening surge). These are the
transition hours where the load curve changes fastest. At hour 8,
the MAPE loss is $-24\%$ --- the official is 24\% more accurate.
This is likely because the official forecast incorporates temperature
and weather information that determines the \emph{timing} of the
ramp; our model uses only historical load patterns.

% ============================================================
\section{ACF Diagnostics: Are the Errors Random?}
% ============================================================

\subsection{What the ACF measures}

After fitting any model, we ask: are the residuals white noise?
If errors follow a systematic pattern, the model has missed something
that could in principle be captured.

The \textbf{autocorrelation function (ACF)} measures:
\[
\rho_k = \text{Corr}(e_t, e_{t-k})
\]
the correlation between today's error and the error $k$ days ago.
If the model is correctly specified, all $\rho_k$ should be near zero.

The \textbf{Ljung-Box test} is a formal test of:
$H_0: \rho_1 = \rho_2 = \cdots = \rho_{14} = 0$ (no autocorrelation
up to lag 14). Rejection means remaining serial structure.

\subsection{What we find}

\textbf{VAR score residuals (in-sample):} fitting the pure VAR(1,7)
without external features on the 2019--2022 training data, the
Ljung-Box statistic is 504 for Score 1 and 553 for Score 2 (both
$p = 0$). This means the scores have remaining autocorrelation
beyond lags 1 and 7 --- which is expected, because holidays and
seasonality create systematic patterns in the scores that the VAR
without features cannot capture. This is the \emph{motivation} for
adding external features.

\textbf{Curve forecast errors (out-of-sample):} for the full main
model with all features, the Ljung-Box statistic at each of the 24
hours ranges from 197 to 7,227, all with $p = 0$. This means
even the best model leaves strongly autocorrelated errors.

\begin{intuitionbox}
\textbf{Why are the out-of-sample errors so autocorrelated?}
Look at the distribution across hours. Hour 6 has LB = 7,227.
This is the morning ramp start. The model systematically
over-forecasts or under-forecasts the ramp timing for consecutive
days --- during Ramadan the ramp shifts to a different pattern and
persists for weeks; during summer the ramp is shallower and this
also persists. These regime patterns create serial dependence in
errors that cannot be removed by lagged features alone.
\end{intuitionbox}

\subsection{Why this matters: the residual step}

The key insight is: \emph{if errors are predictable, we can correct
for them}. The \textbf{residual second step} does exactly this.
For target day $d$, after producing the base forecast $\hat{Y}_d^{base}(h)$:

\begin{enumerate}[nosep]
  \item Compute past fitted residuals: $\hat{e}_s(h) = Y_s(h) - \hat{Y}_s^{base}(h)$
    for all $s < d$ (strictly past data only).
  \item Estimate a correction curve $\hat{r}_d(h)$ from recent past residuals.
  \item Apply the correction: $\hat{Y}_d^{res}(h) = \hat{Y}_d^{base}(h) + \lambda_d \hat{r}_d(h)$
    where $\lambda_d$ is chosen by rolling-origin cross-validation.
\end{enumerate}

This is rolling-origin throughout: no information leakage. The
massive error autocorrelation --- especially at ramp hours --- is
exactly what the residual step leverages. It wins at all 24 of 24
hours including the 5 hours where the base model loses to the official forecast.

% ============================================================
\section{Eigenfunction Stability and Dynamic FPCA}
% ============================================================

A natural question about FPCA is: do the eigenfunctions $\phi_1, \phi_2$
remain stable over time, or do they shift as the Turkish grid evolves?
If they shift, we might need \textbf{dynamic FPCA} --- a more complex
method that allows time-varying eigenfunctions.

We tested this directly:

\begin{enumerate}[nosep]
  \item Fit FPCA separately on three sub-periods: 2019--2021, 2020--2022, 2022--2024.
  \item Compute the \textbf{inner product} between eigenfunctions from
    different periods. An inner product of 1 means identical shape;
    0 means orthogonal (completely rotated).
\end{enumerate}

Results: all pairwise inner products exceed 0.993. The variance
proportion explained by FPC1 stays between 92.4\% and 93.2\%
across all sub-periods. The eigenfunctions are stable.

\begin{resultbox}
\textbf{Conclusion:} dynamic FPCA is not needed for this dataset.
Turkish electricity load shapes --- the overnight trough, morning
ramp, midday plateau, evening peak --- have not structurally changed
over 2019--2024. Static FPCA is appropriate, and we can say so with
empirical evidence rather than just assumption.
\end{resultbox}

% ============================================================
\section{The Random Forest Benchmark}
% ============================================================

\subsection{Why a machine learning comparison?}

The score regression in the main model is linear OLS:
$\xi_d = a + A_1\xi_{d-1} + A_7\xi_{d-7} + Bz_d + u_d$.

A natural question: does the relationship between the features $z_d$
and the scores $\xi_d$ have nonlinear structure that OLS misses?

\textbf{Random forest (RF)} answers this. RF on the same feature
vector $z_d$ plus lagged scores can capture any smooth nonlinear
function of the inputs without specifying the functional form.

\subsection{Why keep the FPCA structure}

We apply RF to the \emph{score regression} only --- we still use FPCA
to represent the curves and reconstruct the forecast curve from
predicted scores. This keeps the functional data analysis framework
intact. The RF is a drop-in replacement for OLS, not an alternative to FPCA.

\subsection{What the comparison tells us}

\begin{itemize}[nosep]
  \item If \textbf{FPCA-OLS $\approx$ FPCA-RF}: the linear score equation
    is adequate. We get the same accuracy with a structured, interpretable
    model. This is the preferred outcome for the thesis's FDA identity.
  \item If \textbf{FPCA-RF $\gg$ FPCA-OLS}: there is meaningful nonlinearity
    in the feature-score relationship. This motivates a future extension but
    does not invalidate the current results.
\end{itemize}

[Results to be added after the RF computation completes.]

% ============================================================
\section{The Coherent Narrative}
% ============================================================

Here is the full story in one flow:

\begin{enumerate}
  \item \textbf{The object:} daily electricity consumption is a curve, not
    a vector. Treating it as a curve preserves intraday shape.

  \item \textbf{The compression:} FPCA reduces 24 correlated observations
    to 2 uncorrelated scores that explain 97.3\% of curve variation.
    The eigenfunctions are empirically stable over 6 years.

  \item \textbf{The dynamics:} scores follow a VAR(1,7) with calendar,
    holiday, load-shape, and derivative features. The features are
    essential: without them, the model is worse than the official forecast.

  \item \textbf{The gain:} the main model (K=2, all features) beats the
    official forecast by 19.2\% MAE and 18.7\% MAPE, with a
    Diebold-Mariano p-value of $10^{-163}$. The improvement is real.

  \item \textbf{Where it loses:} transition hours (6--8, 18--19) where
    the official forecast's weather information gives it an edge. The
    DM test shows this explicitly.

  \item \textbf{The residual step:} ACF diagnostics show that errors are
    strongly autocorrelated, especially at transition hours. The
    residual step exploits this predictable pattern and achieves 46\%
    MAE gain --- winning at all 24 of 24 hours.

  \item \textbf{The benchmark:} random forest on the same features tests
    whether OLS leaves nonlinearity on the table. [Result pending.]

  \item \textbf{Statistical validity:} rolling-origin evaluation ensures
    no leakage. DM tests confirm the gains are statistically significant.
    ACF diagnostics characterise the remaining structure.
\end{enumerate}

% ============================================================
\section{Concepts to Know Cold for the Presentation}
% ============================================================

\begin{tabular}{p{0.25\textwidth}p{0.65\textwidth}}
\toprule
\textbf{Term} & \textbf{One-sentence explanation} \\
\midrule
B-spline & Piecewise polynomial basis for approximating smooth functions \\
FPCA & PCA applied to curves; gives eigenfunctions and daily scores \\
Score & The projection of a curve onto an eigenfunction; a scalar \\
Eigenfunction & The ``principal component'' in function space; a 24-hour curve \\
VAR(1,7) & Multivariate AR with lags 1 and 7; estimated by OLS equation by equation \\
Rolling-origin & Forecasting analogue of cross-validation; train only on strictly past data \\
MAE & Mean absolute error in MWh; primary accuracy metric \\
MAPE & Mean absolute percentage error; scale-free version of MAE \\
RMSE & Root mean squared error; penalises large errors more \\
DM test & Hypothesis test for equal forecast accuracy; $t$-test on loss differential \\
ACF & Autocorrelation function; measures serial dependence in residuals \\
Ljung-Box & Portmanteau test: are ACFs jointly zero up to lag $q$? \\
Residual step & Second-stage correction using past errors; exploits error autocorrelation \\
Expanding window & Training set grows as new data arrives; appropriate when structure is stable \\
Inner product & Cosine similarity between two eigenfunctions; 1 = identical, 0 = orthogonal \\
\bottomrule
\end{tabular}

\end{document}
